% BHU PYQ -- Manually transcribed
% Paper: MTB-201 : Calculus-III
% Exam: B.A. (Hons.) Semester II Examination, 2022-23

\documentclass[11pt,a4paper]{article}
\usepackage[a4paper,top=2.5cm,bottom=2.5cm,left=2.8cm,right=2.8cm,headheight=14pt]{geometry}
\usepackage{amsmath,amssymb}
\usepackage[T1]{fontenc}
\usepackage[utf8]{inputenc}
\usepackage{lmodern,microtype,fancyhdr,enumitem,hyperref,lastpage,parskip}

\hypersetup{colorlinks=true,urlcolor=black,linkcolor=black,
  pdftitle={MTB-201 Calculus-III -- BHU B.A. Sem II 2022-23}}

\pagestyle{fancy}\fancyhf{}
\renewcommand{\headrulewidth}{0.4pt}\renewcommand{\footrulewidth}{0.4pt}
\fancyhead[L]{\small   \textit{Banaras Hindu University}}
\fancyhead[R]{\small   \textit{MTB-201: Calculus-III}}
\fancyfoot[C]{\small Page  \thepage\ of \pageref{LastPage}}


\newlist{parts}{enumerate}{2}
\setlist[parts,1]{label=(\alph*),leftmargin=*,itemsep=5pt,topsep=3pt}
\setlist[parts,2]{label=(\roman*),leftmargin=*,itemsep=3pt,topsep=2pt}

\newcommand{\pts}[1]{\hfill   \textit{[#1]}}


\begin{document}

\begin{center}
  {\large\bfseries BANARAS HINDU UNIVERSITY}\\[2pt]
  {\normalsize B.A.\ (Hons.) --- Semester II Examination, 2022-23}\\[6pt]
  \rule{0.8\linewidth}{0.5pt}\\[6pt]
  {\large\bfseries Paper No.\ MTB-201: Calculus-III}\\[4pt]
  {\normalsize Time: 3 Hours \hfill Full Marks: 70}\\[2pt]
  {\small Ref.\ No.:\ 063898}
\end{center}
\rule{\linewidth}{0.4pt}
\smallskip



\noindent   \textit{Instructions: Answer   \textbf{five} questions including Question~1 which is compulsory. Figures in right-hand margin indicate marks.}
\bigskip




\noindent   \textbf{Question 1.}   \textit{(Compulsory --- 2 marks each.)}

\begin{parts}
  \item Determine the simultaneous limit and repeated limits of $f(x,y)=\dfrac{x^3-y^3}{x^2+y^2}$ at $(0,0)$.
  \item By using the definition of limit, show that $\displaystyle\lim_{(x,y) o(1,2)}(x+y^2)=3$.
  \item If $z=x^{4/3}y^{2/3}$, where $x=t^3$ and $y=t^2$, find $\dfrac{dz}{dt}$.
  \item Investigate for maximum and minimum values of $y=\sin x+\cos 2x$.
  \item If $f(x,y)=\dfrac{xy}{\sqrt{x^2+y^2}}$ for $x^2+y^2
eq0$, find $f_x(0,0)$ and $f_y(0,0)$.
  \item Verify Euler's theorem for $u=x^2\log\!\left(\dfrac{y}{x}\right)$.
  \item Show that $\displaystyle B(n,n)=\frac{\sqrt{\pi}\,\Gamma(n)}{2^{2n-1}\,\Gamma\!\left(n+ frac{1}{2}\right)}$.
  \item If $u^a v^b w^c=1$, show that $\dfrac{\partial^2 w}{\partial u\,\partial v}=\dfrac{-1}{(w\log e w)^3}$ when $u=v=w$.
  \item Evaluate $\displaystyle\int_0^1 x^4(1-x)^3\,dx$.
  \item Show that $\displaystyle\int_0^{\pi/2}\sqrt{\sin \theta}\,d \theta\cdot\int_0^{\pi/2}\frac{d \theta}{\sqrt{\sin \theta}}=\pi$.
  \item Evaluate $\displaystyle\int_0^{\infty}e^{-x^2}dx$.
\end{parts}
\medskip\rule{\linewidth}{0.2pt}

\medskip


\noindent   \textbf{Question 2.}

\begin{parts}
  \item State and prove a necessary and sufficient condition for the existence of limit of $f(x,y)$ at $(a,b)$. Evaluate $\displaystyle\lim_{(x,y) o(0,0)}\frac{x^2y^2}{x^2y^2+(x-y)^2}$.\pts{6}
  \item If $f(x,y)$ possesses partial derivatives at every point of a neighbourhood of $(a,b)$ and these are continuous at $(a,b)$, show that $f$ is differentiable at $(a,b)$.\pts{6}
\end{parts}
\medskip\rule{\linewidth}{0.2pt}

\medskip


\noindent   \textbf{Question 3.}

\begin{parts}
  \item If $u=f(r)$ where $r^2=x^2+y^2$, prove $\dfrac{\partial^2 u}{\partial x^2}+\dfrac{\partial^2 u}{\partial y^2}=f''(r)+\dfrac{1}{r}f'(r)$.\pts{6}
  \item Find the extreme values of $f(x,y)=x^3+y^3-3xy$.\pts{6}
\end{parts}
\medskip\rule{\linewidth}{0.2pt}

\medskip


\noindent   \textbf{Question 4.}

\begin{parts}
  \item Change the order of integration and evaluate $\displaystyle\int_0^a\int_0^{\sqrt{a^2-x^2}}\!f(x,y)\,dy\,dx$.\pts{6}
  \item Evaluate $\displaystyle\iint_R\frac{dA}{(1+x^2+y^2)^2}$ where $R$ is the region $x^2+y^2\leq1$.\pts{6}
\end{parts}
\medskip\rule{\linewidth}{0.2pt}

\medskip


\noindent   \textbf{Question 5.}

\begin{parts}
  \item Prove $\Gamma(n+1)=n\,\Gamma(n)$ and hence $\Gamma(n+1)=n!$ for positive integers.\pts{6}
  \item Evaluate $\displaystyle\int_0^{\pi/2}\sin^5 \theta\cos^6 \theta\,d \theta$ using Beta and Gamma functions.\pts{6}
\end{parts}
\medskip\rule{\linewidth}{0.2pt}

\medskip


\noindent   \textbf{Question 6.}

\begin{parts}
  \item State and prove Dirichlet's theorem on multiple integrals. Evaluate $\displaystyle\iiint x^l y^m z^n\,dx\,dy\,dz$ over the region $x+y+z\leq1$, $x,y,z\geq0$.\pts{6}
  \item Find the volume bounded by $z=x^2+y^2$ and $z=2y$.\pts{6}
\end{parts}
\medskip\rule{\linewidth}{0.2pt}

\medskip


\noindent   \textbf{Question 7.}

\begin{parts}
  \item Evaluate $\displaystyle\int_0^{\infty}\frac{x^{p-1}}{1+x}\,dx=\frac{\pi}{\sin(p\pi)}$ for $0<p<1$.\pts{6}
  \item Evaluate $\displaystyle\int_0^a\int_0^{\sqrt{a^2-x^2}}\!\!\int_0^{\sqrt{a^2-x^2-y^2}}xyz\,dz\,dy\,dx$.\pts{6}
\end{parts}

\vfill\rule{\linewidth}{0.4pt}

\begin{center}\small
\href{https://bhu-exam.vercel.app/}{Click here to visit our website}\\[4pt]
\href{https://me-aryan.vercel.app/}{Click here to visit our contributor}\\[4pt]
\href{https://quashan-me.vercel.app/}{Click here to visit our contributor}
\end{center}
\end{document}
